\documentclass[11pt,a4paper]{article}

% Codificación y lenguaje
\usepackage[utf8]{inputenc}
\usepackage[T1]{fontenc}
\usepackage[spanish,es-nodecimaldot]{babel}

% Presentación
\usepackage[a4paper,margin=2.7cm]{geometry}
\usepackage{parskip}
\usepackage{booktabs}
\usepackage{array}
\usepackage{longtable}
\usepackage{xcolor}
\usepackage{fancyhdr}
\usepackage{lastpage}
\usepackage{amsmath}
\usepackage{listings}
\usepackage[hidelinks]{hyperref}

\definecolor{codebg}{RGB}{247,248,250}
\definecolor{codecomment}{RGB}{70,110,70}
\definecolor{codekeyword}{RGB}{0,70,140}
\definecolor{codestring}{RGB}{150,60,30}

\lstdefinelanguage{VerilogCustom}{
  morekeywords={module,endmodule,input,output,inout,wire,reg,parameter,localparam,
    always,begin,end,case,endcase,default,if,else,assign,signed,integer,for,
    initial,display,finish},
  sensitive=true,
  morecomment=[l]{//},
  morecomment=[s]{/*}{*/},
  morestring=[b]"
}

\lstset{
  language=VerilogCustom,
  basicstyle=\ttfamily\small,
  keywordstyle=\color{codekeyword}\bfseries,
  commentstyle=\color{codecomment}\itshape,
  stringstyle=\color{codestring},
  backgroundcolor=\color{codebg},
  frame=single,
  rulecolor=\color{black!20},
  numbers=left,
  numberstyle=\tiny\color{black!50},
  numbersep=8pt,
  breaklines=true,
  columns=fullflexible,
  keepspaces=true,
  showstringspaces=false,
  tabsize=4
}

\pagestyle{fancy}
\fancyhf{}
\lhead{Arquitectura de Computadoras}
\rhead{TP N.\textsuperscript{o} 1 -- ALU}
\cfoot{\thepage\ de \pageref{LastPage}}
\setlength{\headheight}{14pt}

\newcommand{\verilog}[1]{\texttt{#1}}

\title{\textbf{Diseño e implementación de una ALU combinacional parametrizable}\\[0.4em]
       \large Trabajo Práctico N.\textsuperscript{o} 1}
\author{Javier}
\date{1 de septiembre de 2026}

\begin{document}

\begin{titlepage}
  \hypersetup{pageanchor=false}
  \centering
  \vspace*{1.2cm}
  {\Large Universidad \par}
  \vspace{0.3cm}
  {\large Arquitectura de Computadoras \par}
  \vspace{2.2cm}
  {\LARGE\bfseries Diseño e implementación de una ALU combinacional parametrizable\par}
  \vspace{0.5cm}
  {\large Trabajo Práctico N.\textsuperscript{o} 1\par}
  \vfill
  \begin{tabular}{rl}
    Estudiante: & Javier \\
    Fecha: & 1 de septiembre de 2026
  \end{tabular}
  \vfill
\end{titlepage}

\hypersetup{pageanchor=true}
\clearpage
\pagenumbering{roman}
\tableofcontents
\clearpage
\pagenumbering{arabic}

\begin{abstract}
Este informe presenta el diseño de una unidad aritmético-lógica (ALU) de ocho
bits descrita en Verilog y destinada a su implementación en una FPGA. El
núcleo de la ALU se construye exclusivamente con lógica combinacional: su
salida depende en cada instante de dos operandos y de un código de operación,
sin almacenar estado ni utilizar reloj. Se describen la interfaz parametrizable,
las operaciones implementadas, la interpretación de datos con signo y la
integración con un módulo superior que permite cargar los valores mediante
switches y botones. Finalmente, se explica la estrategia de simulación mediante
testbench y se presentan ejemplos de cálculo que permiten verificar el
funcionamiento esperado.
\end{abstract}

\section{Introducción}

Una unidad aritmético-lógica es un bloque fundamental de un procesador: recibe
operandos, ejecuta una operación determinada por una señal de control y entrega
un resultado. En el presente trabajo se implementa una ALU parametrizable en
Verilog, con operandos y resultado de ocho bits y un código de operación de seis
bits.

El propósito central no es escribir un algoritmo secuencial, sino describir el
hardware que una herramienta de síntesis puede construir en la FPGA. En este
contexto, una expresión como

\begin{lstlisting}[numbers=none]
o_data = i_data_a + i_data_b;
\end{lstlisting}

representa la lógica de un sumador. La selección entre varias operaciones se
implementa mediante lógica combinacional equivalente a un multiplexor.

\section{Objetivos}

Los objetivos del trabajo son:

\begin{itemize}
  \item Diseñar una ALU de ocho bits con operaciones aritméticas, lógicas y de
        desplazamiento.
  \item Describir el núcleo funcional utilizando lógica combinacional en
        Verilog, sin registros ni señales de reloj dentro de la ALU.
  \item Parametrizar el ancho de datos y de operación para favorecer la
        reutilización del módulo.
  \item Integrar el núcleo con una interfaz superior para una placa FPGA,
        permitiendo cargar y conservar los operandos desde switches.
  \item Verificar el comportamiento mediante simulación y relacionar las
        formas de onda con los resultados matemáticos esperados.
\end{itemize}

\section{Arquitectura general}

La solución se divide en dos módulos con responsabilidades diferentes:

\begin{description}
  \item[\verilog{ALU.v}] contiene el núcleo combinacional. Recibe los operandos
        \verilog{i\_data\_a}, \verilog{i\_data\_b} y el código \verilog{i\_op};
        produce \verilog{o\_data}.
  \item[\verilog{ALU\_top.v}] constituye el módulo superior. Conecta la ALU
        con el reloj, el reset, los botones, el bus de switches y los LEDs.
        Incluye registros para conservar las entradas cargadas.
\end{description}

La arquitectura puede resumirse de la siguiente manera:

\begin{center}
\fbox{\begin{minipage}{0.82\textwidth}
\centering
\textbf{Interfaz de la placa}\\[0.2em]
Switches + botones + reloj\\[0.4em]
$\downarrow$\\[0.4em]
\textbf{Registros en \verilog{ALU\_top}} \quad
($A$, $B$, \verilog{Op})\\[0.4em]
$\downarrow$\\[0.4em]
\textbf{ALU combinacional}\quad $\longrightarrow$ \quad
Resultado en LEDs
\end{minipage}}
\end{center}

Esta separación es esencial. La ALU propiamente dicha es combinacional; los
registros del módulo superior no forman parte de su cálculo, sino que permiten
que el usuario cargue $A$, luego $B$ y finalmente la operación sin perder los
valores previamente seleccionados. El resultado no se registra: se actualiza
automáticamente cuando cambia alguna de las entradas almacenadas.

\section{Interfaz y parametrización}

La declaración del módulo permite modificar el tamaño de los buses mediante
parámetros:

\begin{lstlisting}[caption={Interfaz parametrizable del núcleo},label={lst:interface}]
module ALU #(
    parameter NB_OP = 6,
    parameter NB_DATA = 8
)(
    input wire signed [NB_DATA-1:0] i_data_a,
    input wire signed [NB_DATA-1:0] i_data_b,
    input wire [NB_OP-1:0] i_op,
    output reg signed [NB_DATA-1:0] o_data
);
\end{lstlisting}

Con la configuración utilizada, $A$ y $B$ son operandos de ocho bits y
\verilog{i\_op} posee seis bits. Si se instancia el módulo con
\verilog{NB\_DATA = 16}, los operandos y el resultado pasan a ser de dieciséis
bits sin modificar la estructura de la ALU. El calificador \verilog{signed}
permite interpretar los operandos de datos en complemento a dos para las
operaciones aritméticas y para el desplazamiento aritmético.

\section{Operaciones implementadas}

Los códigos de operación se definen mediante \verilog{localparam}, lo que
mejora la legibilidad del circuito y evita repetir constantes sin nombre en la
descripción. La Tabla~\ref{tab:operations} resume el conjunto implementado.

\begin{table}[ht]
\centering
\caption{Códigos y operaciones de la ALU.}
\label{tab:operations}
\renewcommand{\arraystretch}{1.15}
\small
\begin{tabular}{c c c >{\raggedright\arraybackslash}p{6.8cm}}
\toprule
\textbf{Hex.} & \textbf{Binario} & \textbf{Operación} & \textbf{Descripción}\\
\midrule
\texttt{20} & \texttt{100000} & ADD & $A+B$\\
\texttt{22} & \texttt{100010} & SUB & $A-B$\\
\texttt{24} & \texttt{100100} & AND & $A\mathbin{\&}B$\\
\texttt{25} & \texttt{100101} & OR  & $A\mathbin{|}B$\\
\texttt{26} & \texttt{100110} & XOR & $A\mathbin{\hat{~}}B$\\
\texttt{03} & \texttt{000011} & SRA & $A$ desplazado aritméticamente a la derecha $B$ posiciones\\
\texttt{02} & \texttt{000010} & SRL & $A$ desplazado lógicamente a la derecha $B$ posiciones\\
\texttt{27} & \texttt{100111} & NOR & $\mathord{\sim}(A\mathbin{|}B)$\\
\bottomrule
\end{tabular}
\end{table}

En las operaciones de desplazamiento, el segundo operando se utiliza como
cantidad de desplazamiento. La diferencia entre SRA y SRL se observa cuando el
operando $A$ es negativo: SRA replica el bit de signo, mientras que SRL
introduce ceros en las posiciones más significativas.

\section{Implementación de la lógica combinacional}

La lógica del núcleo se describe mediante un bloque \verilog{always @(*)}. La
lista sensible implícita indica que el bloque debe reevaluarse cuando cambia
cualquiera de sus entradas. La asignación inicial garantiza un valor definido
para la salida y evita la inferencia accidental de un latch.

\begin{lstlisting}[caption={Selección combinacional de operaciones},label={lst:logic}]
always @(*) begin
    o_data = {NB_DATA{1'b0}};

    case (i_op)
        OP_ADD: o_data = i_data_a + i_data_b;
        OP_SUB: o_data = i_data_a - i_data_b;
        OP_AND: o_data = i_data_a & i_data_b;
        OP_OR:  o_data = i_data_a | i_data_b;
        OP_XOR: o_data = i_data_a ^ i_data_b;
        OP_NOR: o_data = ~(i_data_a | i_data_b);
        OP_SRL: o_data = i_data_a >> i_data_b;
        OP_SRA: o_data = \$signed(i_data_a) >>> i_data_b;
        default: o_data = {NB_DATA{1'b0}};
    endcase
end
\end{lstlisting}

La sentencia \verilog{case} selecciona la operación asociada con el código de
control. Conceptualmente, cada operación alimenta una entrada de un
multiplexor y \verilog{i\_op} determina cuál de ellas llega a la salida. La
variable \verilog{o\_data} se declara como \verilog{reg} porque recibe
asignaciones dentro de un bloque \verilog{always}; esto no significa que se
implemente un registro físico. Al no existir un flanco de reloj en este bloque,
la síntesis produce lógica combinacional.

Para códigos no contemplados se entrega cero. Además de hacer determinista el
comportamiento, esta decisión completa la asignación de la salida para todos
los caminos posibles.

\section{Integración con la interfaz de la FPGA}

El módulo superior utiliza un bus compartido de switches y tres botones de
carga. En cada flanco ascendente del reloj se almacena el valor presente en el
bus según el botón presionado:

\begin{lstlisting}[caption={Carga de entradas en el módulo superior},label={lst:top}]
always @(posedge clock) begin
    if (reset) begin
        reg_data_A <= {(DATA_LEN){1'b0}};
        reg_data_B <= {(DATA_LEN){1'b0}};
        reg_op     <= {(OP_LEN){1'b0}};
    end
    else if (button_1)
        reg_data_A <= bus[DATA_LEN-1:0];
    else if (button_2)
        reg_data_B <= bus[DATA_LEN-1:0];
    else if (button_3)
        reg_op     <= bus[OP_LEN-1:0];
end
\end{lstlisting}

La secuencia de uso es la siguiente:

\begin{enumerate}
  \item Configurar el valor de $A$ en los switches y presionar el botón 1.
  \item Configurar el valor de $B$ y presionar el botón 2.
  \item Configurar el código de operación en los seis bits inferiores y
        presionar el botón 3.
  \item Observar el resultado en los LEDs.
\end{enumerate}

La asignación no bloqueante \verilog{<=} es apropiada para este bloque
secuencial, ya que representa la actualización simultánea de registros en un
flanco de reloj. El reset síncrono inicializa los tres registros. Luego, la
instancia de la ALU recibe sus valores:

\begin{lstlisting}[caption={Instanciación del núcleo en el módulo superior},label={lst:instance}]
ALU #(
    .NB_OP(OP_LEN),
    .NB_DATA(DATA_LEN)
) alu (
    .i_data_a(reg_data_A),
    .i_data_b(reg_data_B),
    .i_op(reg_op),
    .o_data(led)
);
\end{lstlisting}

Las conexiones por nombre hacen explícita la correspondencia entre cada puerto
y la señal interna, reduciendo el riesgo de errores de conexión.

\section{Ejemplos de funcionamiento}

Para ilustrar la operación, considérese $A=5$ y $B=3$. Las salidas esperadas
para algunas operaciones son:

\begin{table}[ht]
\centering
\caption{Ejemplos de resultados con operandos positivos.}
\label{tab:examples}
\renewcommand{\arraystretch}{1.15}
\begin{tabular}{c c c c}
\toprule
\textbf{Operación} & \textbf{Código} & \textbf{Cálculo} & \textbf{Resultado}\\
\midrule
ADD & \texttt{20} & $5+3$ & $8$\\
SUB & \texttt{22} & $5-3$ & $2$\\
AND & \texttt{24} & \texttt{00000101 \& 00000011} & \texttt{00000001}\\
OR  & \texttt{25} & \texttt{00000101 | 00000011} & \texttt{00000111}\\
XOR & \texttt{26} & \texttt{00000101 \string^ 00000011} & \texttt{00000110}\\
NOR & \texttt{27} & $\mathord{\sim}(5\mathbin{|}3)$ & \texttt{11111000}\\
\bottomrule
\end{tabular}
\end{table}

El resultado de NOR se muestra en ocho bits y equivale a $-8$ en complemento
a dos. Esto evidencia que las operaciones lógicas trabajan sobre el patrón de
bits completo, aunque el bus de datos pueda interpretarse con signo.

Como segundo ejemplo, para ADD se pueden seleccionar $A=2$, $B=\texttt{FF}$ y
el código \texttt{20}. En ocho bits con signo, \texttt{FF} representa $-1$;
por lo tanto,

\begin{equation}
  2 + (-1) = 1,
\end{equation}

y la salida esperada es \texttt{01}. La aritmética se realiza con el ancho del
bus, por lo que los resultados que exceden el rango representable se conservan
en los ocho bits menos significativos.

\section{Verificación mediante simulación}

La verificación funcional se realiza con \verilog{tb\_ALU.v}, un testbench que
instancia directamente el núcleo. Esta elección permite probar la ALU sin
incluir botones, reloj ni restricciones físicas de la placa. El estímulo asigna
operandos, cambia el código de operación, espera un intervalo de tiempo y
observa \verilog{o\_data}.

En una simulación comportamental se deben visualizar, como mínimo, las
siguientes señales:

\begin{center}
\begin{tabular}{c c c c}
\texttt{i\_data\_a[7:0]} &
\texttt{i\_data\_b[7:0]} &
\texttt{i\_op[5:0]} &
\texttt{o\_data[7:0]}
\end{tabular}
\end{center}

Para facilitar la lectura, los datos pueden observarse
en decimal con signo o hexadecimal, mientras que el código de operación se
puede observar en binario o hexadecimal.

El testbench cambia la operación cada $10\,\mathrm{ns}$ y conserva el mismo par
de operandos durante la secuencia. Los intervalos relativos son:

\begin{table}[ht]
\centering
\caption{Secuencia temporal de operaciones del testbench.}
\label{tab:timing}
\begin{tabular}{c c}
\toprule
\textbf{Intervalo} & \textbf{Operación}\\
\midrule
0--10 ns   & ADD\\
10--20 ns  & SUB\\
20--30 ns  & AND\\
30--40 ns  & OR\\
40--50 ns  & XOR\\
50--60 ns  & SRA\\
60--70 ns  & SRL\\
70--80 ns  & NOR\\
\bottomrule
\end{tabular}
\end{table}

Para validar una operación se debe ubicar el cursor de la forma de onda en el
interior de un intervalo estable, evitando una transición, y comparar el valor
de salida con la ecuación correspondiente. Por ejemplo, si se observan
\verilog{i\_data\_a = 02}, \verilog{i\_data\_b = FF}, \verilog{i\_op = 20} y
\verilog{o\_data = 01}, la suma verificada es $2+(-1)=1$.

El módulo superior también puede simularse con \verilog{tb\_ALU\_TOP.v}, pero
requiere generar correctamente el reloj, el reset, el bus y los botones. Por
ello, para verificar primero la función de la ALU se recomienda utilizar
\verilog{Test\_Bench\_ALU}; \verilog{ALU\_top} representa la integración con
la placa y no constituye por sí mismo un estímulo de simulación.

\section{Consideraciones de síntesis}

Desde el punto de vista de hardware, el núcleo combina:

\begin{itemize}
  \item operadores aritméticos para suma y resta;
  \item compuertas bit a bit para AND, OR, XOR y NOR;
  \item lógica de desplazamiento para SRA y SRL;
  \item lógica de selección para entregar al puerto de salida la operación
        elegida.
\end{itemize}

No hay memoria interna en \verilog{ALU.v}. El tiempo de respuesta queda
determinado por el retardo de propagación de la red combinacional sintetizada.
En cambio, los registros de \verilog{ALU\_top.v} son necesarios para desacoplar
la entrada manual de los switches del cálculo: una vez cargados, mantienen los
operandos aunque el bus cambie.

El archivo de restricciones \texttt{.xdc} completa la implementación física al
asociar los puertos del módulo superior con los pines de la placa FPGA. De esta
manera, la descripción lógica, la interfaz de usuario y la asignación física
quedan separadas en archivos con responsabilidades claras.

\section{Conclusiones}

Se diseñó una ALU parametrizable de ocho bits capaz de realizar ocho
operaciones aritméticas, lógicas y de desplazamiento. La implementación de
\verilog{ALU.v} cumple el objetivo principal del trabajo: el resultado se
obtiene exclusivamente a partir de las entradas actuales y no depende de un
reloj ni conserva estado.

La integración mediante \verilog{ALU\_top.v} mostró por qué una aplicación
física necesita registros aun cuando el núcleo sea combinacional. Estos
registros almacenan $A$, $B$ y \verilog{Op}, mientras que la ALU calcula y
actualiza el resultado de forma inmediata ante cualquier cambio de dichos
valores. La división modular facilita tanto la reutilización del núcleo como
su verificación independiente mediante testbench.

Finalmente, la simulación permite relacionar cada código de operación con la
ecuación y el patrón de bits observados en la forma de onda. Este procedimiento
constituye una verificación funcional previa a la síntesis y a la prueba sobre
la placa FPGA.

\end{document}
